\documentclass[11pt]{article}
\usepackage[margin=1in]{geometry}
\usepackage{amsmath,amssymb,amsthm}
\title{Disproof of Conjecture 00000004236}
\author{earthking11}
\date{October 6, 2026}
\newtheorem{proposition}{Proposition}
\begin{document}
\maketitle
\section*{Counterexample}
Let $G$ be the triangle on vertices $0,1,2$ with nonnegative weights
\[
 w_{01}=1,\qquad w_{02}=1,\qquad w_{12}=3.
\]
Root a shortest-path tree at $0$.  The edges $01$ and $02$ form such a tree:
their direct path weights are $1$, whereas the alternative paths have weights
$w_{02}+w_{21}=4$ and $w_{01}+w_{12}=4$, respectively.

After the two tree edges are removed, the only residual edge is $12$, of
weight $3$.  Thus the minimum residual edge weight is $3$.

There is exactly one nontrivial simple $1$-cycle in a triangle.  Its weight is
\[
 w_{01}+w_{12}+w_{20}=1+3+1=5.
\]
Therefore the optimal nontrivial cycle has weight $5$, not $3$.

\begin{proposition}
The weight of an optimal $1$-cycle need not equal the minimum residual edge
weight after deleting a shortest-path tree, even for nonnegative weights.
\end{proposition}
\begin{proof}
The triangle above has optimal cycle weight $5$ and residual minimum $3$.
Since $5\ne3$, it is a counterexample.
\end{proof}

The missing term is structural: a non-tree edge determines a fundamental
cycle only after the tree path between its endpoints is added.  Its bare edge
weight is therefore not, in general, the weight of that cycle.

The Lean project checks every numerical and finite combinatorial assertion.
\end{document}
